\section{Introduction}\label{sec:introduction}

A Hecke eigensheaf realizes a local system on a curve by the geometry of
modifications of bundles. Over a finite field, an arithmetic local system
contains more information than its geometric restriction: it also specifies
Frobenius. The corresponding automorphic object must therefore carry a
Frobenius structure that preserves its eigenisomorphisms. We construct such
objects with Borel level and tame regular-unipotent monodromy, after a finite
extension of the field of constants.

\subsection{The result}

For a smooth projective pointed curve $(X,x)$, let $G=\SL_n$ and
let $B$ be its upper-triangular Borel. The stack
$\Bun_{G,B,x}(X)$ parametrizes $G$-bundles on $X$ with a $B$-reduction
at $x$. A sheaf on this stack is locally constructible perverse if its
pullback to each smooth finite-type scheme chart, shifted by the relative
dimension of that chart, is bounded constructible and perverse. Its
parabolic global nilpotent cone $\Lambda_{\mathrm{par}}$ consists of
generically nilpotent Higgs fields with at most a simple pole at $x$ and
residue in the nilradical of the chosen Borel. Singular support is
interpreted on the same smooth charts.

For representations $\boldsymbol V=(V_i)_{i\in I}$ of the dual group
$\PGL_n$, write $\Hecke_{I,\boldsymbol V}$ for the Hecke functor with
legs in $(X\setminus\{x\})^I$. We use the relative Satake normalization
of Section~\ref{sec:conventions}, with no moving-leg shift. In particular,
the tensor-unit functor is exterior product with the constant
$\E$-sheaf on $(X\setminus\{x\})^I$.

A Weil structure over $\F_{q^m}$ means an isomorphism between the
$q^m$-Frobenius pullback of the geometric sheaf and the sheaf itself.
We require the eigenisomorphisms to respect this structure and the
Frobenius structure on the parameter.

\begin{theorem}\label{thm:main}
Let $n\geq2$, let $\F_q$ have characteristic $p>n$, let $\ell\ne p$,
and put $E=\E$. Let $X_0/\F_q$ be a smooth projective geometrically
connected curve of genus at least two, with $x_0\in X_0(\F_q)$, and
put $U_0=X_0\setminus\{x_0\}$. Suppose
\[
 \rho:\pi_1^{\mathrm{et}}(U_0)\longrightarrow\PGL_n(L)
\]
is continuous for a finite extension $L/\Q_\ell$, its geometric
restriction is Zariski dense, and its inertia at $x_0$ kills wild inertia
and has the form
\begin{equation}\label{eq:regular-monodromy}
 \rho(\gamma)=\exp\bigl(t_\ell(\gamma)N\bigr)
\end{equation}
for a regular nilpotent $N$. Here a choice of compatible
$\ell$-power roots of unity identifies the tame quotient
$\Z_\ell(1)$ with $\Z_\ell$ and defines $t_\ell$.

There are an integer $m\geq1$ and a nonzero $E$-adic Weil sheaf $M_0$ on
\[
 \Bun_{\SL_n,B,x_0}(X_0)_{\F_{q^m}}
\]
whose geometric pullback $M$ is locally constructible perverse and has
singular support in $\Lambda_{\mathrm{par}}$. If $\rho_m$ denotes the
restriction of $\rho$ to $\pi_1^{\mathrm{et}}(U_{0,\F_{q^m}})$, then
for every finite set $I$ and every family
$V_i\in\Rep_E(\PGL_n)$ there are isomorphisms of Weil sheaves
\begin{equation}\label{eq:main-eigenmaps}
 \Hecke_{I,\boldsymbol V}(M_0)
   \simeq M_0\boxtimes\bigboxtimes_{i\in I}(V_i)_{\rho_m}.
\end{equation}
They are natural in the representations and coherent for the unit,
convolution, permutations, and all fusion diagonals.
\end{theorem}

Write $k=\overline{\F}_q$, $X=X_0\times_{\F_q}k$, $x=(x_0)_k$, and
$U=U_0\times_{\F_q}k$ for the geometric curve and its punctured complement.
The Frobenius structure is required on the entire eigenstructure, rather
than only on the underlying perverse sheaf. The finite extension in the
theorem arises when a point on an auxiliary bundle stack is made invariant
under a power of Frobenius. We use Weil structures, so a compatible action
of the infinite cyclic group is sufficient.

\subsection{Background and the geometric input}

The geometric Langlands program replaces automorphic functions by sheaves
on moduli stacks of bundles and asks that Hecke eigenvalues be realized as
local systems on a curve. Drinfeld's rank-two work and Laumon's geometric
formulation established central parts of this perspective
\cite{Drinfeld,Laumon}. Frenkel, Gaitsgory, and Vilonen constructed
unramified $\GL_n$ eigensheaves, including the finite-field Weil setting;
Gaitsgory subsequently gave a geometric proof of the vanishing theorem
underlying that construction \cite{FGV,GaitsgoryVanishing}. Thus the
arithmetic compatibility of an eigenstructure is a classical requirement,
not a new condition introduced here.

For general groups, the restricted spectral stack and its action on
automorphic sheaves with nilpotent singular support were developed by
Arinkin, Gaitsgory, Kazhdan, Raskin, Rozenblyum, and Varshavsky
\cite{AGKRRV}. Gaitsgory and Raskin prove the restricted geometric
Langlands equivalence in characteristic zero and construct its
positive-characteristic specialization; in positive characteristic their
general-group result concerns a union of spectral components
\cite{GR}. We use the characteristic-zero equivalence and the spectral
action, with their precise scopes, in Section~\ref{sec:unramified}.

Our geometric input is the companion manuscript
\emph{Constructible tame Hecke eigensheaves in positive characteristic}
\cite{CT}. It constructs geometric, locally constructible perverse
eigensheaves with Borel level, parabolic nilpotent singular support, and
the full fusion-compatible eigenstructure. Its method first creates level
structure by a characteristic-zero degeneration to a separating node and
then specializes to positive characteristic. We use its nearby-cycle,
nonvanishing, perverse-truncation, and enhancement-descent theorems, stated
at their uses below. Those results do not supply Frobenius descent.
The present paper proves the additional arithmetic compatibility for
$\SL_n$ under the hypotheses of Theorem~\ref{thm:main}.

In characteristic-zero de Rham sheaf theory, F\ae rgeman constructs
coherent tame eigenobjects for irreducible regular-singular parameters,
conditional on the spectral-action and localization-compatibility inputs
specified in \cite[Sections~1.3.6--1.4.3]{Faergeman}. Regular holonomicity
and generic perversity in the tame setting remain expectations in that
preprint \cite[Remark~1.4.1.2]{Faergeman}. Our theorem establishes
perversity in the geometric $\ell$-adic setting and equips the
eigenstructure with a finite-field Weil action, for the dense parameters
with regular-unipotent boundary monodromy specified above.

The nodal geometry is part of the automorphic gluing framework of Nadler
and Yun \cite{NYAutomorphicGluing}, with twisted curves as in
Abramovich--Vistoli and Olsson \cite{AbramovichVistoli,OlssonTwisted}.
Root stacks and proper henselian invariance
of finite covers allow us to retain finite-quotient monodromy through the
two degenerations \cite{LieblichOlsson,Rydh}. The sheaf-theoretic
specialization used here is geometric nearby cycles, in the form developed
by Hansen--Scholze and Gaitsgory--Raskin and adapted to level structure
in \cite{CT,HansenScholze,GR}.

\subsection{The arithmetic constructions and the proof}

There are two extra difficulties in making the geometric construction
equivariant. First, an invariant parameter need not make an arbitrarily
chosen geometric eigenobject visibly invariant. Second, the two boundary
parameters at a separating node must be glued by a conjugation that
preserves both Frobenius and compact image.

Section~\ref{sec:unramified} treats the first difficulty. For a dense
unramified $\PGL_n$-parameter in characteristic zero, the relevant
restricted spectral component is a smooth formal space with one point.
Transport of the spectral action preserves that component. The inverse
image of its skyscraper has scalar endomorphisms and no negative
self-extensions; these properties recognize its transport up to shift.
An invariant polarization on the curve rules out a nonzero shift on the
automorphic side, giving a generator isomorphism on this particular
eigenobject. Density and the
trivial center of $\PGL_n$ force compatibility with its full eigenstructure.
This recognition argument avoids choosing an equivariant normalization
of the Langlands equivalence itself.

Section~\ref{sec:parameters} first lifts the original arithmetic parameter
to a punctured curve $C_1$ in characteristic zero. At its puncture the
Frobenius matrix $A$ and the regular nilpotent logarithm $N_1$ satisfy
$\Ad(A)N_1=qN_1$. A cocharacter commuting with $A$ scales $N_1$.
We choose a ramified cover $C_2\to C_1$ whose degree makes the required
gluing conjugator lie in a fixed compact open subgroup. The resulting
boundary identification commutes with $A$. This compact equivariant gluing
is the second reusable ingredient: inverse boundary monodromies alone
would not ensure a continuous compact-valued parameter on the smoothing.

The compatible finite torsor towers on the two punctured components glue
on balanced twisted nodal models and lift to a dense unramified parameter
on their smooth generic curve. Apply the equivariant unramified result
there. In Section~\ref{sec:nodal}, nearby cycles on a prescribed node type,
restriction to one component, and descent from enhanced to ordinary flags
produce an equivariant perverse eigensheaf on $C_1$. The restriction uses
a point defined over a finite extension of the characteristic-zero base;
a power of the generator fixes that point. Finally,
Section~\ref{sec:final} specializes this eigensheaf to the original
finite-field curve. Both specializations use a separate support-closure
argument to prove nonvanishing. Naturality of the comparisons in
Section~\ref{sec:conventions} carries the entire eigenstructure through
these steps.
